\chapter{The shift register, read as a deforming surface}

\textbf{Purpose:} Record what the deformation of a harmonic reading of the SHA-256 compression state is, prove it belongs to the state and not to the observer or the probe, and trace it to the one structural feature of the compression function that could produce it. \textbf{Scope:} \texttt{examples/\allowbreak{}proofing/\allowbreak{}state\_deflection.py}, and the figures in \texttt{docs/\allowbreak{}figures/\allowbreak{}state-deformation/}.

A harmonic reading of the SHA-256 compression state renders a surface, and that surface deforms as the round clock turns. The observer-invariance making the reading trustworthy is stated first, because it is the calibration letting the rest be read as measurement and not as viewpoint.

\section{The object}

SHA-256 compresses in sixty four rounds over a state of eight thirty-two bit words, which is \(256\) bits. The reading places those \(256\) bits on \(256\) points of a ring arrangement and expands the lit set in spherical harmonics. The state becomes a field on a sphere, and the field has a shape. Bit \(i\) lights point \(i\); no mapping was invented to make the instrument fit the object, because \(256\) bits and \(256\) points are the same count.

The compression agrees with the published digest of the empty string before any of this is read. The object under the instrument is SHA-256 and not an approximation of it.

\section{The calibration: which part of the shape is the observer's}

A shape read this way has two anisotropies, and they answer to a rotation of the observer differently. This is standard representation theory of the rotation group, stated here because it is the null the measurement rests on, and not because it is new.

\textbf{Deflection}, the power per degree \(P_l = \sum_m |a_{lm}|^2\), is invariant under every rotation in \(\mathrm{SO}(3)\): the degree-\(l\) subspace carries a unitary irreducible representation of the group, and a magnitude built from it cannot record how the object was turned. Measured null under a whole-step rotation: \(\mathbf{2.539 \times 10^{-14}}\), the arithmetic and not the object.

\textbf{Torsion}, the phase \(\arg a_{lm}\) of the same coefficients, moves by exactly \(-m\alpha\) under a rotation by \(\alpha\). It records the turn, sign included. Measured recovery of a known angle: \(\mathbf{1.803 \times 10^{-13}}\) radians.

The magnitude of the anisotropy is therefore observer-independent and its direction is observer-dependent, and the split is exact. This is a property of the reading and holds for any lit set, SHA-256 or otherwise. Its role here is precise: a change in the shape's \textbf{magnitude} across the clock is a change in the object, while a change in its \textbf{direction} could be either the object or the frame. Every claim below is made on a magnitude, where the observer has been divided out, or is checked against a control holding the frame fixed.

\section{The measurement, on the surface itself}

The reading renders exactly one deformed surface in the scene, and it is the state's. Every enclosing shell measures a vertex-radius standard deviation of exactly zero, and each is a perfect sphere. The state's surface measures a standard deviation of \(0.999\) against a mean radius of \(8.795\), at \(11.4\) percent. Read as a radius field over the sphere and tracked across the sixty four rounds:

\begin{center}
\begin{tabular}{@{}llp{0.42\linewidth}@{}}
\toprule
feature & degree & across the clock \\
\midrule
the elongation, a lemon & 2 & \(q_{zz} \approx -1.50\) at every round, \(17\) percent of the mean radius \\
the offset, a cone & 1 & \(d_z\) swings about \(0.07\) and changes sign several times \\
\bottomrule
\end{tabular}
\end{center}

The quadrupole is a fixed feature: the surface is elongated by the same amount, in the same axis, the whole way through. The dipole migrates and inverts. The cone walks around the surface and changes which pole it favors more than once.

\section{The claim, stated as a number}

The dipole does not merely move. It moves \textbf{coherently}: its sign, taken round by round, comes in runs, and not in the alternation independent signs would give. Twenty runs against an expected \(30.5\), at a run-length test assuming nothing about the values themselves:

\begin{center}
\(z = -2.87\)
\end{center}

Measured a second way, on a different object through different code, the rendered mesh in the viewer against the lit set in Python, the same statistic reads \(\mathbf{z = -3.0}\). Two pipelines with no shared arithmetic land two hundredths apart.

A drift at that level is what a surface computed from a slowly changing state looks like. The next section is the control pinning the change to the state.

\section{Pinning the change to the state}

Two interventions, because the round clock advances everything at once and an observation that the picture changes when the clock moves does not by itself say what moved it.

\textbf{The probe was moved and nothing happened.} Holding the round fixed and swinging the beam \(180\) degrees in azimuth, \(128\) in elevation and \(4.7\) times in radius left the surface unchanged in every digit: mean radius, dipole and quadrupole all identical before, during and after. One round of the clock, with the beam untouched, then moved the dipole by five times. The surface answers to the state and not to the probe.

\textbf{The rounds were replaced by independent states.} The real round sequence gives \(z = -2.87\). A control swapping the rounds for independent states of the same per-round weights, everything else identical, gives \(z = 0.00\) with a spread of \(0.94\). The real sequence sits \(3.1\) control deviations out. The drift is in the succession of states, and not in the reading of any one of them.

\section{The mechanism, which is not exotic}

SHA-256's round computes two words and shifts the other six:

\begin{center}
\begin{tabular}{@{}ll@{}}
\toprule
\(a' = T_1 + T_2\) & \(e' = d + T_1\) \\
\(b' = a\) & \(f' = e\) \\
\(c' = b\) & \(g' = f\) \\
\(d' = c\) & \(h' = g\) \\
\bottomrule
\end{tabular}
\end{center}

Six of the eight words at round \(r\) are round \(r-1\)'s words in a new position, measured at \(63\) of \(63\) rounds. A surface computed from a state that is five-sixths carried over cannot jump between consecutive rounds, and a coherent, autocorrelated drift is therefore the shift register showing.

\textbf{That is the claim entire, and it is worth stating without deflation.} The reading was built to show a field. Nobody built it to expose the Merkle-Damgard shift register, and the register turned up in it anyway: legible to the eye as a lemon whose cone walks and inverts, and holding at three standard deviations when the drift is measured against a control. The result is not a new property of SHA-256. It is a demonstration that a harmonic boundary reading makes a known structural property of the compression function visible and measurable, with the observer divided out and the cause established by intervention.

\section{Scope, stated plainly}

\textbf{The observer-invariance is classical.} Rotation-invariant power and orientation-carrying phase are standard \(\mathrm{SO}(3)\) representation theory, the same fact behind an angular power spectrum being usable without its phases. It is the calibration here, and not a finding.

\textbf{What is specific to SHA-256} is the behavior of the anisotropy across the clock: a persistent quadrupole, a dipole inverting in runs at \(z = -2.87\), and the trace to the shift register by intervention.

\textbf{No claim is made} about a weakness in SHA-256, about its output, or about anything a cryptographic reduction would touch. The reading is taken on the intermediate compression state, which carries the shift structure by design. The digest, being the feed-forward of all sixty four rounds, is a different object and is not read here.

\section{Reproducing it}

\begin{center}
\begin{tabular}{@{}ll@{}}
\toprule
command & what it returns \\
\midrule
\texttt{state\_deflection.py --check} & the instrument, against the digest \\
\texttt{state\_deflection.py --shape} & dipole and quadrupole per round \\
\texttt{state\_deflection.py --runs} & the autocorrelation against a control \\
\bottomrule
\end{tabular}
\end{center}

The figures, and the intervention table, are in \texttt{docs/\allowbreak{}figures/\allowbreak{}state-deformation/}.
